\subsection{问题一：需求、时长与资源配置模型}

\subsubsection{需求预测与峰平判定}

对来源$s$、日期$d$，记开单事件数为$N_{sd}$。为检验星期与趋势结构，比较三个不使用留出期信息的基准模型：7日季节朴素、最近8周同星期中位数和日历趋势岭回归。后者写为
\begin{equation}
  \widehat N_{sd}=\beta_0+\beta_1 t_d+\sum_{w=1}^{6}\gamma_w I(w_d=w)
  +\sum_{m=2}^{12}\eta_m I(m_d=m)+\rho_1\sin\frac{2\pi t_d}{365.25}
  +\rho_2\cos\frac{2\pi t_d}{365.25}.
\end{equation}
岭参数只在训练期设计矩阵上确定。采用
\begin{equation}
 \mathrm{WAPE}=\frac{\sum_d|N_{sd}-\widehat N_{sd}|}{\sum_dN_{sd}}
\end{equation}
作为主评价，因为体检低需求日会使逐日百分比误差不稳定。留出期结果表明：住院和门诊开单量均由最近8周同星期中位数取得最低WAPE，分别为16.04\%和14.04\%；体检开单量由日历趋势模型取得最低WAPE，但仍为69.83\%，说明团检尖峰不能靠常规日历变量充分预测。三类合计开单量的最低WAPE为24.92\%。医院日需求预测中显式建模星期结构并滚动验证是必要步骤\cite{luo2017,lin2019}，但体检业务还需要接入未来团检计划，不能只依赖历史时间序列。

对留出期每日需求分布，以中位数附近$[Q_{0.45},Q_{0.55}]$定义平峰，以$Q_{0.90}$定义高峰，以$Q_{0.95}$定义压力日。住院开单日平峰为197人，高峰阈值283人；门诊分别为156和198人；体检因团检尖峰，平峰48人而$Q_{0.90}$达到308.2人。故设备配置采用“固定能力底座+高峰弹性块”，而不是按体检全年均值永久占用机房。

\subsubsection{服务时长的分层估计}

设项目$j$属于类别$k(j)$，题给时长区间为$[l_k,u_k]$，名义时长
\begin{equation}
  p_j^{(0)}=5\left\lceil\frac{(l_k+u_k)/2}{5}\right\rceil,
  \qquad
  p_j^{(1)}=5\left\lceil\frac{u_k}{5}\right\rceil.
\end{equation}
5\%特殊事件使用$p_j^{(1)}$，其余使用$p_j^{(0)}$。各项目族参数见表\ref{tab:duration}。心脏、产科III/IV级的绝对占用显著高于一般项目，是高峰配置中必须保护的长任务。

同一申请含多个项目时，主排程按项目逐项累计上述时长；在同一兼容机房连续完成可以省去转运，但不把不同项目视为同时扫描。训练期另以0.8、0.9和1.0三种合并系数检查需求量口径，本文不把缺少现场证据的节约比例写入患者级可行域。

\begin{table}[H]
  \centering
  \caption{项目族调度时长及训练期间隔证据}
  \label{tab:duration}
  \begin{tabular}{lrrrr}
\toprule
项目族 & 下界/min & 名义/min & 上界/min & 训练期代理样本数 \\
\midrule
产科III/IV级 & 30 & 35 & 40 & 13,451 \\
产科I/II级 & 15 & 17.5 & 20 & 20,095 \\
心脏 & 10 & 12.5 & 15 & 141,017 \\
介入/定位 & 5 & 7.5 & 10 & 6,782 \\
血管 & 5 & 7.5 & 10 & 59,006 \\
经阴道/腔内 & 5 & 7.5 & 10 & 21,244 \\
泌尿/经腹妇科 & 5 & 7.5 & 10 & 75,137 \\
腹部/胃肠 & 5 & 7.5 & 10 & 99,719 \\
浅表器官 & 5 & 7.5 & 10 & 67,825 \\
儿科专项 & 5 & 7.5 & 10 & 3,950 \\
一般其他 & 5 & 7.5 & 10 & 3,400 \\
\bottomrule
\end{tabular}

\end{table}

为描述医生异质性，先选取同一医生、同一天的单项目事件，计算相邻报告间隔$g_n$，仅保留$0<g_n\leq60$分钟。对医生$i$，由有效工作日吞吐量得到对数相对因子$h_i$，再作经验贝叶斯收缩：
\begin{equation}
  \widetilde h_i=\omega_i h_i+(1-\omega_i)\mu_h,
  \qquad
  \omega_i=\frac{m_i}{m_i+\lambda},
\end{equation}
其中$m_i$为有效工作日数，$\mu_h$为总体对数吞吐均值。医生—项目组合的调度时长取直接间隔中位数与题给名义时长经医生因子修正后的收缩组合，并截断在$[l_k,u_k]$内。该方法形成42\,336个完整组合、6\,416个直接证据组合。由于相邻报告间隔包含非扫查活动，$\widetilde h_i$只用于相对速度与能力分析，患者级主排程仍采用题给绝对时长，以免把代理噪声固化为硬约束。基于检查类别预测超声检查时长具有现实依据\cite{lai2025}，但本文对字段缺失作了更保守的解释。

\subsubsection{项目—设备能力证据}

设备能力从表1的检查项目文本逐项建立。对项目名称$q$和设备$r$定义三级匹配：
\begin{equation}
  A_{qr}=\begin{cases}
  2,&\text{项目名称精确、包含或语义核心匹配；}\\
  1,&\text{无项目级证据但项目族有具体能力证据；}\\
  3,&\text{文本明确为床旁且$r=7$；}\\
  0,&\text{否则。}
  \end{cases}
\end{equation}
主情景允许$A_{qr}\geq1$，同时记录$A_{qr}=1$的回退任务；严格敏感性只允许$A_{qr}\in\{2,3\}$。宽泛描述如“常规项目”不能单独建立某个项目族的硬能力。图\ref{fig:coverage}显示，住院、门诊任务的严格或床旁证据分别达到98.2\%和91.2\%，体检只有69.4\%，其余30.6\%依赖类别回退。这一差异决定了问题三不能只报告总兼容率。

\begin{figure}[H]
  \centering
  \includegraphics[width=0.92\textwidth]{fig03_capability_coverage.pdf}
  \caption{三类患者项目—设备兼容证据覆盖}
  \label{fig:coverage}
\end{figure}

关键资源见表\ref{tab:capability}. 床旁仅7号机可执行；产科III/IV级集中于64、43、65号三台设备；儿科专项仅45、44号两台。心脏、血管等项目虽有较多可用设备，但长时长和高需求仍会形成负荷瓶颈。

\begin{table}[H]
  \centering
  \caption{关键资源的设备覆盖与结构判断}
  \label{tab:capability}
  \begin{tabularx}{\textwidth}{l r X l}
\toprule
关键资源 & 可用设备数 & 设备编号 & 结构判断 \\
\midrule
床旁 & 1 & 7 & 唯一瓶颈 \\
心脏 & 5 & 73|113|48|45|49 & 多房间 \\
血管 & 12 & 40|69|72|71|41|42|46|47|44|74|75|77 & 多房间 \\
III/IV产科 & 3 & 64|43|65 & 少数房间 \\
I/II产科 & 13 & 40|41|42|46|47|44|64|43|65|160|74|75|77 & 多房间 \\
儿科专项 & 2 & 45|44 & 少数房间 \\
介入/定位 & 15 & 40|69|72|41|42|46|47|44|64|43|65|160|74|75|77 & 多房间 \\
\bottomrule
\end{tabularx}

\end{table}

\subsubsection{医生共享时隙容量}

表6提供医生每日项目数及首末活动时刻。对训练日$d$和5分钟槽$t$，若$t$位于医生$i$重构活动区间内，则$I_{idt}=1$。星期$w$的主容量和偏紧容量分别为
\begin{equation}
  c^{(0)}_{wt}=\min\left\{22,\operatorname{Med}_{d:\,w(d)=w}\sum_iI_{idt}\right\},
  \quad
  c^{(L)}_{wt}=\min\left\{22,Q_{0.25,d:\,w(d)=w}\sum_iI_{idt}\right\}.
\end{equation}
首活动在7--9时之间归整为8时开班，12:30--14时归整为13时开班；末活动在11--12:30归整为12时，在16--19时归整为17时，其他情况保留观察区间并向前扩10分钟。该规则只修复报告时间代理对班次边界的轻微偏移，不把单个医生复制成全天可用。

\begin{figure}[H]
  \centering
  \includegraphics[width=\textwidth]{fig02_doctor_capacity.pdf}
  \caption{训练期重构的星期—时隙医生共享容量}
  \label{fig:doctor-capacity}
\end{figure}

\begin{table}[H]
  \centering
  \caption{各星期医生并发容量范围}
  \label{tab:doctor-capacity}
  \begin{tabular}{lrrrrrr}
\toprule
 & \multicolumn{3}{c}{主情景} & \multicolumn{3}{c}{偏紧情景} \\
\cmidrule(lr){2-4}\cmidrule(lr){5-7}
星期 & 最小 & 中位 & 最大 & 最小 & 中位 & 最大 \\
\midrule
周一 & 15 & 17.5 & 20 & 13 & 15 & 17 \\
周二 & 15 & 17 & 20 & 14 & 15.5 & 17 \\
周三 & 15 & 17 & 20 & 13 & 15.5 & 17 \\
周四 & 16 & 18 & 20 & 14 & 15.5 & 18 \\
周五 & 15 & 16.5 & 19 & 13 & 14.5 & 17 \\
周六 & 6 & 6 & 6 & 5 & 5 & 5 \\
周日 & 5 & 5 & 5 & 4 & 4 & 5 \\
\bottomrule
\end{tabular}

\end{table}

工作日主容量为15--20人，周六为6人，周日为5人；实际可同时运行的设备数取设备数22与$c_{wt}$的较小值。图\ref{fig:doctor-capacity}表明午后容量普遍略低，周末更明显。医生下四分位情景将用于检验人力不足风险。

\subsubsection{峰平配置方案}

综合需求、能力和医生容量，提出以下配置原则。

\begin{enumerate}
  \item 7号机全天保留床旁任务，非床旁项目只有在确定不会挤占床旁窗口时才可回填；这是唯一确定性硬瓶颈。
  \item 64、43、65号设备形成产科III/IV级保护池；45、44号形成儿科保护池。其余兼容机房按项目稀缺度分配，灵活项目优先使用覆盖面较广且当前负荷较低的设备。
  \item 平峰按训练期中位需求开放固定能力底座；高峰不固定新增某台“万能机”，而是在上午空腹块和下午非空腹块分别释放弹性时隙。体检团检高峰需提前导入团检名单，单独预留成块容量。
  \item 5\%特殊病例不另设固定机房，而在每个工作块保留与上界时长差对应的缓冲。医生并发低于设备数时，先补医生班次而不是继续开放设备。
\end{enumerate}

表\ref{tab:peak-flat-plan}把上述原则落实为22台现有设备的运行配置。这里的“开放”表示允许排程占用，并不等于临时购置设备；任一时刻真正可并行工作的台数仍受$c_{wt}$限制。

\begin{table}[H]
  \centering
  \caption{22台现有设备的峰平运行配置}
  \label{tab:peak-flat-plan}
  \begin{tabularx}{\textwidth}{p{0.19\textwidth} c X X}
    \toprule
    资源池 & 台数 & 平峰配置 & 高峰/压力配置 \\
    \midrule
    床旁7号机 & 1 & 床旁任务保护，确认无待服务床旁任务时才回填 & 保持床旁专用，不以普通任务填满 \\
    产科III/IV级（64、43、65） & 3 & 保留专项能力，按实际产科项目启用 & 三台均进入专项候选，空余槽才承接兼容任务 \\
    儿科专项（45、44） & 2 & 保留儿科能力，空余槽可按证据回填 & 两台均保护儿科窗口，避免一般项目提前占用 \\
    其余兼容设备 & 16 & 按$Q_{0.45}$--$Q_{0.55}$工作量滚动开放 & 按$Q_{0.90}$需求和已知团检名单释放成块时隙；$Q_{0.95}$压力日先增医生班次 \\
    \bottomrule
  \end{tabularx}
\end{table}

该方案既回答了“哪些机器配置哪些项目”，又保留了患者级模型根据当日项目组合动态选择具体机房的空间；固定配置负责保护稀缺能力，滚动排程负责利用可替代能力。
